\section{Cofibered categories, bifibered categories}
\label{VI.10}
Let us consider a category $\cal{F}$ over~$\cal{E}$,
with the projection functor
$$
p\colon \cal{F}\to \cal{E},
$$
it defines a category $\cal{F}^\circ$ over
$\cal{E}^\circ$, by the projection functor
$$
p^\circ\colon \cal{F}^\circ\to \cal{E}^\circ.
$$
A morphism $\alpha\colon \eta\to\xi$ in $\cal{F}$ is said to be
\emph{cocartesian} if it is a cartesian morphism for
$\cal{F}^\circ$ over~$\cal{E}^\circ$. Making this explicit, one sees that this
means that for every object $\xi'$ of~$\cal{F}_S$, the map
$u\mto u\circ \alpha$
$$
\Hom_S(\xi,\xi')\to\Hom_f(\eta,\xi')
$$
is bijective. One then also says that $(\xi,\alpha)$ is a
\emph{direct image}
of~$\eta$ by $f$, in the category $\cal{F}$ over~$\cal{E}$. If
it exists for every $\eta$ in $\cal{F}_T$, one says that the
direct image functor by $f$ exists, and one denotes this functor
% original p. 182
$f_*^{\cal{F}}$ or~$f_*$,
once it has been chosen. It is thus defined by
an isomorphism of bifunctors on~$\cal{F}_T^\circ\times\cal{F}_S$:
$$
\Hom_S(f_*(\eta),\xi)\isomto \Hom_f(\eta,\xi).
$$
Thus if $f_*$ exists, for $f^*$ to exist, it is necessary and sufficient that
$f_*$
admit an adjoint functor, \ie that there exist a functor
$f^*\colon \cal{F}_S\to\cal{F}_T$ and an isomorphism of bifunctors
$$
\Hom_S(f_*(\eta),\xi)\isomto \Hom_T(\eta,f^*(\xi)).
$$
Let $g\colon U\to T$ be another morphism in $\cal{E}$, and suppose
that the inverse images and direct images by $f,g$ and $fg$
exist. Let us consider then the functorial homomorphisms
$$
\begin{array}{cccc}
c^{f,g}\colon&f_*g_*&\dpl\from&(fg)_*\\
c_{f,g}\colon&g^*f^*&\dpl\to&(fg)^*.
\end{array}
$$
One observes that if one regards $f_*g_*$ and $g^*f^*$ as a
pair of adjoint functors, as well as $(fg)_*$ and $(fg)^*$, the two
preceding homomorphisms are adjoint to one another. Thus
one is an isomorphism if and only if the other is. In
particular:
\begin{proposition}
\label{VI.10.1} Suppose that the category $\cal{F}$ over~$\cal{E}$
is prefibered and coprefibered. For it to be
fibered, it is necessary and sufficient that it be
cofibered.
\end{proposition}

Of course, one says that $\cal{F}$ is coprefibered \resp cofibered
over~$\cal{E}$, if $\cal{F}^\circ$ is prefibered \resp fibered over~$\cal{E}$. We shall say that $\cal{F}$ is bifibered
over~$\cal{E}$,
if it is at once fibered and cofibered over~$\cal{E}$.

\section{Various examples}
\label{VI.11}

\begin{enumerate}

\item[a)] \textbf{Categories of arrows of}~$\cal{E}$. Let
$\cal{E}$ be a category. Let us denote by $\mathbf{\Delta^1}$ the
category associated to the totally ordered set
with two elements $[0,1]$;
% original p. 183
it thus has two objects 0 and 1, and besides the two identity
morphisms an arrow $(0,1)$ with source 0 and target 1. Let
$$
\CatFl(\cal{E})=\SheafHom(\mathbf{\Delta^1},\cal{E})
$$
one calls it the \emph{category of arrows of} $\cal{E}$.
The object 1 of~$\mathbf{\Delta^1}$ defines a canonical functor,
called the \emph{target functor}
$$
\CatFl(\cal{E})\to \cal{E}
$$
(the functor defined by the object 0 of~$\mathbf{\Delta^1}$ is
called the \emph{source functor}). For every object $S$ of~$\cal{E}$,
the fiber-category $\CatFl(\cal{E})_S$ is canonically isomorphic
to the category $\cal{E}_{/S}$ of objects of~$\cal{E}$
over~$S$.

Let us consider a morphism $f\colon T\to S$ in $\cal{E}$, then there
corresponds to it a canonical functor
$$
f_*\colon \cal{E}_{/T}=\cal{F}_T\to \cal{E}_{/S}=\cal{F}_S
$$
and a functorial isomorphism
$$
\Hom_S(f_*(\eta),\xi)\isomto\Hom_f(\eta,\xi)
$$
which thus makes of~$f_*$ a direct image functor for $f$ in
$\cal{F}$. One has moreover here
$$
(\id_S)_*=\id_{\cal{F}_S},\quad (fg)_*=f_*g_*,\quad
c^{f,g}=\id_{(fg)},
$$
\ie $\cal{F}$ is equipped with a co-splitting over~$\cal{E}$. A fortiori,
$\cal{F}$ is cofibered over~$\cal{E}$. Let us now note that
the set of morphisms in $\cal{F}$ is in bijective
correspondence with the set of commutative square diagrams in
$\cal{E}$.
$$
\xymatrix{
X\ar[d]_-u&\ar[l]_-{f'}Y\ar[d]^-v\\
S&\ar[l]_-{f}T
}
$$
By
% original p. 184
definition, the morphism in question is cartesian if the
square is cartesian in $\cal{E}$, \ie if it makes of~$Y$ a
fibered product of~$X$ and $T$ over~$S$. The inverse image functor
$f^*$
thus exists if and only if for every object $X$ over~$S$, the fibered
product $X\times_ST$ exists. It follows from~\ref{VI.10.1} that if
the product of two objects over a third always exists in~$\cal{E}$, \ie if $\cal{F}$ is prefibered over~$\cal{E}$,
then $\cal{F}$ is even fibered over~$\cal{E}$.

\medskip
\item[b)] \textbf{Category of presheaves or sheaves on
variable spaces}

Let $\cal{E}=\mathbf{Top}$ be the category of topological spaces.
If $T$ is a topological space, we shall denote by $\cal{U}(T)$ the
category of open sets of~$T$, where the morphisms are the
inclusion maps. If $\cal{C}$ is a category, a
functor $\cal{U}(T)^\circ\to \cal{C}$ is called a
\emph{presheaf} on~$T$ with values in $\cal{C}$, and a
\emph{sheaf} if it satisfies a left exactness condition
which we do not repeat here. The \emph{category
$\cal{P}(T)$ of presheaves on~$T$ with values in
$\cal{C}$} is by definition the category
$\SheafHom(\cal{U}(T)^\circ,\cal{C})$, and the category
$\cal{F}(T)$ of sheaves on~$T$ with values in $\cal{C}$ is the
full subcategory whose objects are the objects of
$\SheafHom(\cal{U}(T)^\circ,\cal{C})$ which are sheaves. If
$f\colon T\to S$ is a morphism in $\cal{E}$, \ie a
continuous map of topological spaces, there corresponds to it by
the increasing map $U\mto f^{-1}(U)$ a functor
$\cal{U}(S)\to \cal{U}(T)$, whence a functor
$$
f_*\colon\SheafHom(\cal{U}(T)^\circ,\cal{C})\to
\SheafHom(\cal{U}(S)^\circ,\cal{C})
$$
called the \emph{direct image functor of presheaves by}
$f$. One sees immediately that the direct image of a sheaf is a
sheaf, thus the functor
$f_*\colon \cal{P}(T)\to\cal{P}(S)$
induces a functor, likewise denoted
$f_*\colon\cal{F}(T)\to\cal{F}(S)$.
One moreover verifies trivially (by the associativity of the
composition of functors) that one has, for a second continuous
map $g\colon U\to T$, the identity
$$
(gf)_* =g_*f_*,\qquad
\text{likewise }\ (\id_S)_*=\id_{\cal{P}(S)}.
$$
In this way, one has obtained a functor
$$
S\mto \cal{P}(S)
$$
\resp $$
S\mto \cal{F}(S)
$$
from
% original p. 185
$\cal{E}$ to~$\Cat$. In fact, we are interested in the
corresponding functor
$$
S\mto \cal{P}(S)^\circ,\qquad\text{\resp~}S\mto
\cal{F}(S)^\circ.
$$
It defines a cofibered category, and even a
co-split one, over the category of topological spaces, which one
calls the \emph{cofibered category of presheaves}
(\resp \emph{sheaves}) \emph{with values in} $\cal{C}$
(understood: on variable spaces). Making explicit the construction
of~no.~\ref{VI.8}, one sees that a morphism from a presheaf $B$ on~$T$
to a presheaf $A$ on~$S$ is a pair $(f,u)$ consisting
of a continuous map from~$T$ to $S$, and of a morphism $u\colon
A\to f_*(B)$ in the category $\cal{P}(S)$. This description
is equally valid for morphisms of sheaves, $\cal{F}$
being a full subcategory of~$\cal{P}$.

In the most important cases, the category $\cal{P}$ and the
category $\cal{F}$ over~$\cal{E}$ are also
fibered categories, \ie for every continuous map, the
direct image functors $\cal{P}(T)\to \cal{P}(S)$ and
$\cal{F}(T)\to\cal{F}(S)$ have an adjoint functor, which is then
denoted $f^*$ and called the inverse image functor of
presheaves \resp the inverse image functor of sheaves, by
the continuous map $f$. This functor exists for example if
$\cal{C}=\Ens$. One can show that the functor $f^*\colon
\cal{P}(S)\to \cal{P}(T)$ exists whenever in $\cal{C}$ the
inductive limits (relative to diagrams in the Universe
under consideration) exist. The question is less easy for
$\cal{F}$; indeed, one notes that (even in the case
$\cal{C}=\Ens$) the inverse image of a presheaf which is a
sheaf is in general not a sheaf, in other words the
inverse image functor of sheaves is not isomorphic to the functor
induced by the inverse image functor of presheaves (despite
the common notation $f^*$). Thus, $\cal{F}$ is a
cofibered subcategory of~$\cal{P}$, but not a
fibered subcategory, \ie \emph{the inclusion functor
$\cal{F}\to\cal{P}$ is not fibrant}.

The cofibered category $\cal{P}$ can be deduced from a more
general cofibered (or rather fibered) category,
obtained as follows. For every category $\cal{U}$ (in
the fixed Universe), one sets
$$
\cal{P}(\cal{U})=\SheafHom(\cal{U},\cal{C})
$$
and
% original p. 186
one notes that $\cal{U}\mto\cal{P}(\cal{U})$ is in a
natural way a contravariant functor in~$\cal{U}$, from the category
$\Cat$ to $\Cat$. It thus defines a split category
over $\cal{E}=\Cat$, which we shall denote $\Cat_{/\!/\cal{C}}$. The
objects of this category are the pairs $(\cal{U},p)$ of a
category $\cal{U}$ and a functor $p\colon\cal{U}\to\cal{C}$,
and a morphism from~$(\cal{U},p)$ to $(\cal{V},q)$ is
essentially a pair $(f,u)$, where $f$ is a functor
$\cal{U}\to\cal{V}$ and $u$ a homomorphism of functors $u\colon p\to
qf$. We leave to the reader the task of making explicit the composition of
morphisms in $\Cat_{/\!/\cal{C}}$. The projection functor
$$
\cal{F}=\Cat_{/\!/\cal{C}}\to\cal{E}=\Cat
$$
associates to the pair $(\cal{U},p)$ the object $\cal{U}$; the
fiber-category at $\cal{U}$ is the category
$\SheafHom(\cal{U},\cal{C})$ (up to isomorphism). When
inductive limits exist in $\cal{C}$, one shows easily
that the fibered category $\Cat_{/\!/\cal{C}}$ over~$\Cat$ is
likewise cofibered over~$\Cat$, \ie one can define the
notion of \emph{direct image of a functor} $p\colon\cal{U}\to\cal{C}$
by a functor $f\colon\cal{U}\to\cal{V}$. The category of
presheaves is deduced from the preceding fibered category
by the change of base
$$
\mathbf{Top}^\circ\to\Cat
$$
(functor $S\mto\cal{U}(S)$ defined above), which gives a
fibered category over~$\mathbf{Top}^\circ$, and passing to
the opposite category, one obtains the
cofibered category $\cal{P}$ of presheaves over
$\mathbf{Top}$. The notion of inverse image of a functor corresponds
to that of direct image of a presheaf, the notion of direct
image of a functor to that of inverse image of a
presheaf.

\medskip
\item[c)] \textbf{Objects with operators over an object
with operators}

Let $\cal{F}$ be a category over~$\cal{E}$, and let $S$ be an object of
$\cal{E}$ on which a group $G$ operates, on the left to fix
ideas. This object with operators can be interpreted as
corresponding to a functor $\lambda\colon\cal{E}'\to \cal{E}$ from
the category (with a single object, having $G$ as group
of endomorphisms) $\cal{E}'$ defined by $G$, into the
category $\cal{E}$, and thus there is defined by change of base
a category $\cal{F}'$ over~$\cal{E}'$, which is
fibered \resp cofibered when $\cal{F}$ is so
over~$\cal{E}$. A
% original p. 187
section of~$\cal{E}'$ over~$\cal{F}'$ (necessarily
cartesian, since $\cal{E}'$ is a groupoid, and every
isomorphism in $\cal{F}'$ is cartesian by virtue
of~\ref{VI.6.12}), can also be interpreted as an
$\cal{E}$-functor $\cal{E}'\to\cal{F}$ over~$\lambda$, or
also as an object with operators $\xi$ in $\cal{F}$
``over'' the object with operators $S$.

\medskip
\item[d)] \textbf{Pairs of quasi-inverse adjoint functors;
autodualities}

When the base-category $\cal{E}$ is reduced to two
objects $a,b$ and, besides the identity arrows, to two
isomorphisms $f\colon a\to b$ and $g\colon b\to a$ inverse to one
another (\ie $\cal{E}$ is a connected rigid groupoid with two
objects), a normalized cloven category over~$\cal{E}$ is
essentially the same thing as the system formed by
two categories $\cal{F}_a$ and $\cal{F}_b$ and a \emph{pair of
adjoint functors} $G\colon \cal{F}_a\to\cal{F}_b$ and $F\colon
\cal{F}_b\to\cal{F}_a$, which are equivalences of
categories (hence quasi-inverse to one another). One will take for
$\cal{F}_a$ and $\cal{F}_b$ the fiber-categories of~$\cal{F}$,
for $F$ and $G$ the functors $f^*$ and $g^*$, and the two
isomorphisms
$$
u\colon FG\isomto\id_{\cal{F}_a}\qquad v\colon
GF\isomto\id_{\cal{F}_b}
$$
are $c_{g,f}$ and $c_{f,g}$. The two usual
compatibility conditions between $u$ and $v$ are none other than the condition
\ref{VI.7.4}~$B)$ for the composites $fgf$ and $gfg$. It is easy
to show that these conditions suffice to imply that one indeed has
a pseudofunctor $\cal{E}^\circ\to\Cat$.

An interesting case is that where one has
$$
\cal{F}_b=\cal{F}_a^\circ,\quad G=F^\circ,\quad v=u^\circ.
$$
One calls \emph{autoduality}
in a category $\cal{C}$, the
data of a functor $D\colon \cal{C}\to\cal{C}^\circ$, and of an
isomorphism $u\colon DD^\circ\isomto\id_{\cal{C}}$, such that $u$ and
the isomorphism $u^\circ\colon D^\circ D\isomto\id_{\cal{C}^\circ}$
make of~$(D,D^\circ)$ a pair of adjoint functors,
(necessarily quasi-inverse to one another). This condition
is written:
$$
D(u(x))=u(D(x))\qquad \text{for every $x\in\Ob(\cal{C})$.}
$$
\item[e)]
% original p. 188
\textbf{Categories over a discrete category $\cal{E}$.} One says that $\cal{E}$ is a
\emph{discrete category} if every arrow therein is an
identity arrow, so that $\cal{E}$ is defined up to
unique isomorphism by the knowledge of the set
$I=\Ob(\cal{E})$. The data of a category $\cal{F}$
over~$\cal{E}$ is thus equivalent (up to unique isomorphism)
to the data of a family of categories
$\cal{F}_i$ ($i\in I$), the fiber-categories. Every category
$\cal{F}$ over~$\cal{E}$ is fibered, every $\cal{E}$-functor
$\cal{F}\to\cal{G}$ is cartesian, one has a canonical isomorphism
$$ \SheafHom_{\cal{E}/-}(\cal{F},\cal{G}) \isomto
\prod_{i}\SheafHom(\cal{F}_i,\cal{G}_i).
$$
In particular, one obtains
$$
\mathbf{\Gamma}(\cal{F}/\cal{E})=\varprojLim\cal{F}/\cal{E}\isomto\prod_{i} \cal{E}_i.
$$
\item[f)] \textbf{Suppose that $\cal{E}$ has exactly two objects
$S$ and $T$,} and besides the \textbf{identity morphisms, a
morphism }$f\colon T\to S$. Then a category $\cal{F}$
over~$\cal{E}$ is defined, up to unique $\cal{E}$-isomorphism,
by the data of two categories $\cal{F}_S$
and $\cal{F}_T$ and of a bifunctor $H(\eta,\xi)$ on~$\cal{F}_T^\circ\times\cal{F}_S$, with values in $\Ens$. Indeed,
if~$\cal{F}$ is a category over~$\cal{E}$, one associates to it
the two fiber-categories $\cal{F}_S$ and~$\cal{F}_T$ and
the bifunctor $H(\eta,\xi)=\Hom_f(\eta,\xi)$. One leaves to the reader the
task of making explicit the construction in the opposite direction. For the
category under consideration to be fibered (or prefibered,
which amounts to the same) it is necessary and sufficient that the functor $H$ be
representable with respect to the argument $\xi$. For it
to be cofibered, it is necessary and sufficient that $H$ be
representable with respect to the argument~$\eta$.
\item[g)] Let $\cal{F}=\cal{C}\times\cal{E}$, considered
as a category over~$\cal{E}$ by means of
$\pr_2$. Then $\cal{F}$ is fibered and cofibered over~$\cal{E}$,
and is even equipped with a canonical splitting and co-splitting,
corresponding to the constant functor on~$\cal{E}$, \resp on~$\cal{E}^\circ$, with values in $\Cat$, of value $\cal{C}$. One has
$$
\mathbf{\Gamma}(\cal{F}/\cal{E})\simeq\SheafHom(\cal{E},\cal{C})
$$
and
% original p. 189
$\varprojLim\cal{F}/\cal{E}$ corresponds to the
full subcategory formed of the functors
$F\colon\cal{E\to\cal{C}}$ transforming arbitrary morphisms into
isomorphisms.
\end{enumerate}

\section{Functors on a cloven category}
\label{VI.12}

Let $\cal{F}$ be a normalized cloven category
over~$\cal{E}$. For every object $S$ of~$\cal{E}$ one denotes by
$$
i_S\colon \cal{F}_S\to\cal{F}
$$
the inclusion functor. One thus has a functorial homomorphism, for
every morphism $f\colon T\to S$ in~$\cal{E}$:
$$
\alpha_f\colon i_T f^*\to i_S,
$$
where $f^*$ is the change of base functor
$\cal{F}_S\to\cal{F}_T$ for $f$ defined by the cleavage. Let
now
$$
F\colon\cal{F}\to\cal{C}
$$
be a functor from~$\cal{F}$ to a category $\cal{C}$; set,
for every~$S\in\Ob(\cal{E})$,
$$
F_S=F\circ i_S\colon\cal{F}_S\to\cal{C}
$$
and for every~$f\colon T\to S$ in~$\cal{E}$,
$$
\varphi_f=F*\alpha_f\colon F_T f^*\to F_S
$$
One has thus, to every functor $F\colon \cal{F}\to\cal{C}$,
associated a family $(F_S)$ of functors $\cal{F}_S\to\cal{C}$, and
a family $(\varphi_f)$ of homomorphisms of functors $F_T f^*\to
F_S$. These families satisfy the following conditions:

a) $\varphi_{\id_S}=\id_{F_S}$.

b) For two morphisms $f\colon T\to S$ and $g\colon U\to T$
in~$\cal{E}$, one has commutativity in the square
of functorial homomorphisms:
$$
\xymatrix@C=2cm{ F_U g^* f^* \ar[r]^-{F_U*c_{f,g}} \ar[d]_{\varphi_g * f^*} &
F_U(fg)^* \ar[d]^{\varphi_{fg}} \\ F_T f^*\ar[r]^-{\varphi_f} & F_S.}
$$
% original p. 190
The first relation is trivial, and the second relation
is obtained by applying the functor~$F$ to the commutative diagram
$$
\xymatrix@C=2cm{ g^* f^* (\xi) \ar[r]^-{c_{f,g}(\xi)}
\ar[d]_{\alpha_g(f^*(\xi))} & (fg)^*(\xi) \ar[d]^{\alpha_{fg}(\xi)}\\
f^*(\xi) \ar[r]^-{\alpha_f(\xi)}& \xi }
$$
for a variable object $\xi$ in~$\cal{F}_S$.

If $G$ is a second functor $\cal{F}\to\cal{C}$, giving
rise to functors $G_S\colon\cal{F}_S\to\cal{C}$ and
functorial homomorphisms $\psi_f\colon G_Tf\to G_S$, and if $u\colon
F\to G$ is a functorial homomorphism, then there correspond to it
functorial homomorphisms $u*i_S$:
$$
u_S\colon F_S\to G_S
$$
and one observes immediately that for every morphism $f\colon T\to S$
in $\cal{E}$, one has commutativity in the squares
$$
\begin{array}{c}
\xymatrix@C=1.5cm{ F_Tf^* \ar[r]^-{\phi_f} \ar[d]_{u_T * f^*} & F_S
\ar[d]^{u_S} \\ G_T f^* \ar[r]^-{\psi_f}& G_S}
\end{array}
\leqno\quad \textup{c)}
$$

\begin{proposition}
\label{VI.12.1}
Let $\cal{H}(\cal{F},\cal{C})$ be the category whose objects are
the pairs of families $(F_S)$ ($S\in\Ob(\cal{F})$) of functors
$\cal{F}_S\to\cal{C}$, and of families $(\varphi_f)$
($f\in\Fl(\cal{F})$) of functorial homomorphisms $F_Tf^*\to F_S$,
satisfying the conditions \emph{a) and~b)}, and where the morphisms
are the families $(u_S)$ ($S\in\Ob(\cal{F})$) of homomorphisms
% original p. 191
$F_S\to G_S$, verifying the commutativity condition
\emph{c)} written above, (the composition of morphisms being
effected by the composition of the homomorphisms of functors
$\cal{F}_S\to\cal{C}$). Then the two laws made explicit above
define an \emph{isomorphism} $K$ of the category
$\SheafHom(\cal{F},\cal{C})$ with the category
$\cal{H}(\cal{F},\cal{C})$.
\end{proposition}

It is trivial that one has indeed here a \emph{functor} from the
first category to the second. This functor is fully
faithful, since for given $F,G$, $\Hom(F,G)\to\allowbreak\Hom(K(F),K(G))$
is trivially injective; to show that it is surjective, it suffices
to note that the commutativity condition c) expresses the
functoriality of the maps $u(\xi)=u_S(\xi)\colon
F(S)=F_S(\xi)\to G(\xi)=G_S(\xi)$ for the homomorphisms of the form
$\alpha_f(\xi)$ in $\cal{F}$, on the other hand one has functoriality
on each fiber-category \ie for the morphisms in $\cal{F}$
which are $T$-morphisms ($T\in\Ob(\cal{E})$), whence
functoriality for every morphism in $\cal{F}$,
since an
$f$-morphism (where $f\colon T\to S$ is a morphism in
$\cal{E}$
is in a unique way) a composite of a morphism
$\alpha_f(\xi)$ and a $T$-morphism. It remains therefore to prove that
the functor~$K$ is bijective on objects. The preceding
argument already shows that $K$ is injective on
objects; it remains to prove that it is surjective, \ie that if one starts
from a system $(F_S),(\varphi_f)$, satisfying a) and b); and if one
defines a map $\Ob\cal{F}\to\Ob\cal{C}$ by
$$
F(\xi)=F_S(\xi) \quad\text{for }
\xi\in\Ob\cal{F}_S\subset\Ob\cal{F}
$$
and a map $\Fl(\cal{F})\to\Fl(\cal{C})$ by
$$
F(\alpha_f(\xi)u')=\varphi_f(\xi)\,F_T(u')
$$
for every morphism $f\colon T\to S$ in $\cal{E}$, every object $\xi$
of~$\cal{F}_S$ and every $T$-morphism $u'$ with target $f^*(\xi)$, then one
obtains a \emph{functor} $F$ from~$\cal{F}$ to $\cal{C}$. Indeed,
the relation $F(\id_\xi)=\id_{F(\xi)}$ is trivial, it remains to
prove the multiplicativity $F(uv)=F(u)F(v)$ when one has an
$f$-morphism $u\colon\eta\to\xi$ and a $g$-morphism
$v\colon\zeta\to\nu$, with $f\colon T\to S$ and $g\colon U\to T$
morphisms of~$\cal{E}$. Setting $w=uv$, one will have
$$
u=\alpha_f(\xi)u'\quoi,\quad v=\alpha_g(\eta)v'\quoi,\quad w=\alpha_{fg}(\xi)w'
$$
% original p. 192
with
$$
w'=c_{f,g}(\xi)g^*(u')v'\quad \hbox{(\cf no.~\ref{VI.8})}.
$$

With these notations, one must prove the commutativity of the outer
contour of the diagram below:
$$
\xymatrix{ F_U(\zeta) \ar@{-<} `u[r] `[rrrrrr]^{F(w')} [rrrrrr]
\ar[rr]^{F_u(v')} \ar[drr]_{F(v)} & & F_Ug^*(\eta)
\ar[rr]^{F_ug^*(u')} \ar[d]^{\varphi_g(\eta)} & & F_Ug^*f^*(\xi)
\ar[rr]^{F_U(c_{f,g}(\xi))}\ar[d]^{\varphi_g(f^*(\xi))}& &
F_U(fg)^*(\xi) \ar[d]^{\varphi_{fg}(\xi)} \\ & & F_T(\eta) \ar@{-<}
`d[r] `[rrrr]_{F(u)} [rrrr] \ar[rr]_{F_T(u')} & & F_Tf^*(\xi)
\ar[rr]_{\varphi_f(\xi)} & & F_S(\xi) }
$$
Now the left triangle is commutative by definition of~$F(v)$, the
middle square is commutative since it is deduced from the homomorphism
$u'\colon \xi\to f(\eta)$ by the functorial homomorphism $\alpha_g$,
finally the right-hand square is commutative by virtue of condition
b). The desired conclusion follows.

Suppose now that $\cal{C}$ is likewise a
normalized cloven category over~$\cal{E}$, which we
shall henceforth call $\cal{G}$, and that we are interested
in the $\cal{E}$-functors from~$\cal{F}$ to $\cal{G}$. If $F$ is such a
functor, it induces functors
$$
F_S\colon \cal{F}_S\to\cal{G}_S
$$
for the fiber-categories. On the other hand, for every morphism
$f\colon T\to S$ in~$\cal{E}$, and every object $\xi$ in~$\cal{F}_S$,
the $f$-morphism $F(\alpha_f(\xi))$ factors in a unique way
through a $T$-morphism
$$
\varphi_f(\xi)\colon F_T(f^*_{\cal{F}}(\xi))\to
f^*_{\cal{G}}(F_S(\xi))
$$
(where the $\cal{F}$ or the $\cal{G}$ in the index indicates the
cloven category for which one takes the inverse image
functor), whence a functorial homomorphism of functors from
$\cal{F}_S$ to $\cal{G}_T$:
$$
\varphi_f\colon F_T f^*_{\cal{F}}\to f^*_{\cal{G}}F_S.
$$
% original p. 193
The two systems $(F_S)$ and $(\varphi_f)$ satisfy the
following conditions:

a') $\varphi_{\id_S}=\id_{F_S}$.

b') For two morphisms $f\colon T\to S$ and $g\colon U\to T$ in
$\cal{E}$, one has commutativity in the following diagram of functorial
homomorphisms:
$$
\xymatrix{ F_Ug_F^*f_F^* \ar[rr]^{F_U *c_{f,g}^\cal{F}}
\ar[d]_{\varphi_{g^*}f^*_{\cal{F}}} & & F_U(fg)^*_F
\ar[dd]^{\varphi_{fg}} \\ g^*_{\cal{G}}F_Tf^*_F
\ar[d]_{g^*_{\cal{G}}*\varphi_f} & & \\ g^*_{\cal{G}}f^*_{\cal{G}}F_S
\ar[rr]^{c_{f,g}^\cal{G}*F_S} & & (fg)^*_{\cal{G}}F_S. }
$$
We leave the verification to the reader, as well as
the statement and the proof of the analogue of
proposition~\ref{VI.12.1}, implying that one thus obtains a
bijective correspondence between the set of $\cal{E}$-functors from
$\cal{F}$ to $\cal{G}$, and the set of systems
$(F_S)$,$(\varphi_f)$ satisfying the conditions a') and b')
above. Of course, in this correspondence, the cartesian
functors are characterized by the property that
the homomorphisms $\varphi_f$ are isomorphisms.

\begin{remark}
Of course, it is most often advantageous to reason
directly on fibered categories without using
explicit cleavages, which in particular dispenses with having to appeal,
for the simple notion of~$\cal{E}$-functor or of cartesian $\cal{E}$-functor,
to a cumbersome interpretation as
above. It is to avoid unbearable heaviness, and to
obtain more intrinsic statements,
% original p. 194
that we have had to renounce starting as in \cite{VI.2} from the
notion of cloven category (called ``fibered
category'' in \loccit), which recedes to second rank in favor of
that of fibered category. It is moreover probable that,
contrary to the usage still predominant now,
tied to old habits of thought, it will eventually
prove more convenient in universal problems not to
put the emphasis on \emph{one} solution supposed chosen once
and for all, but to put all solutions on an equal
footing.
\end{remark}

\begin{thebibliography}{9}
\bibitem[1]{VI.1} A\ptbl \textsc{Grothendieck}, Sur quelques points d'alg\`ebre
homologique,
T\^ohoku Math.\ J. \textbf{9} (1957) p\ptbl 119--221.


\bibitem[2]{VI.2} A\ptbl \textsc{Grothendieck}, Technique de descente et
Th\'eor\`emes d'existence,~I, S\'eminaire Bourbaki~190,
December~1959.
\end{thebibliography}

